\section{Conclusion}
\label{sec:conclusion}

In this paper, we introduced \textbf{VulnRAG}, an agentic retrieval-augmented generation framework designed to overcome the ``Retrieval Gap'' in autonomous vulnerability intelligence. By transitioning from static, linear retrieval to a stateful, iterative agentic loop, VulnRAG provides the structural discovery and self-correcting reasoning required for high-stakes security assessments. Our primary contributions—structural-semantic hybrid retrieval via relationship graphs, zero-shot agentic self-correction, Multi-Factor Semantic Reranking (MFSR), and complexity-aware multi-tier routing—collectively redefine the state-of-the-art for vulnerability-specific RAG.

Empirical evaluation on a curated benchmark demonstrates that VulnRAG significantly outperforms traditional vector-based baselines, achieving a 97.4\% reasoning success rate while maintaining an 84.2\% reduction in operational inference costs. Most importantly, our structural discovery mechanism yielded a 2.1$\times$ improvement in identifying operationally linked vulnerabilities, proving that structural context is a critical, yet previously underutilized, signal in the cybersecurity intelligence ecosystem. As the field of autonomous red teaming continues to mature, VulnRAG serves as a robust and economically viable foundation for the next generation of intelligent, self-optimizing security agents.
